\documentclass{article}
\usepackage{amsmath}
\usepackage{listings}
\usepackage{xcolor}
\usepackage{hyperref}

\title{Utility Functions for Environment Variable Access in C++}
\author{}
\date{}

\lstset{
  language=C++,
  basicstyle=\ttfamily\small,
  keywordstyle=\color{blue},
  commentstyle=\color{gray},
  stringstyle=\color{teal},
  showstringspaces=false,
  breaklines=true,
  frame=single
}

\begin{document}

\maketitle

\section*{Overview}

These C++ utility functions provide safer and more expressive alternatives to the standard \lstinline|getenv()| function for retrieving environment variables. Each function handles the case where the environment variable is not set in a distinct way: exiting the program, returning a default value, or throwing an exception.

\section*{Function Descriptions}

\subsection*{1. \lstinline|const char* getenv_or_exit(const char* var_name, void (*exit_fn)(int) = std::exit)|}

This function attempts to retrieve the environment variable \lstinline|var_name| using \lstinline|getenv()|. If the variable is not set (i.e., \lstinline|getenv()| returns \lstinline|nullptr), it calls the specified \lstinline|exit\_fn| function, which defaults to \lstinline|std::exit|.

\textbf{Example Usage:}
\begin{lstlisting}
const char* path = getenv_or_exit("MY_ENV_VAR");
\end{lstlisting}

\textbf{Notes:}
\begin{itemize}
  \item Useful for critical variables that must be present for the program to continue.
  \item Allows custom exit behavior by passing a different exit function.
\end{itemize}

\subsection*{2. \lstinline|const char* getenv_or_default(const char* var_name, const char* default_value)|}

Returns the value of the environment variable \lstinline|var_name| if it exists. Otherwise, returns \lstinline|default_value|.

\textbf{Example Usage:}
\begin{lstlisting}
const char* mode = getenv_or_default("RUN_MODE", "debug");
\end{lstlisting}

\textbf{Notes:}
\begin{itemize}
  \item Useful when a fallback or default behavior is acceptable.
  \item Avoids the need for manual null checks.
\end{itemize}

\subsection*{3. \lstinline|const char* getenv_or_throw(const char* var_name)|}

Returns the environment variable value if it exists. If not, throws a \lstinline|std::runtime_error| with a descriptive message.

\textbf{Example Usage:}
\begin{lstlisting}
try {
    const char* value = getenv_or_throw("API_KEY");
} catch (const std::runtime_error& e) {
    std::cerr << "Missing env var: " << e.what() << std::endl;
}
\end{lstlisting}

\textbf{Notes:}
\begin{itemize}
  \item Best used in exception-enabled code where you want to handle missing variables as recoverable errors.
\end{itemize}

\section*{Conclusion}

These utility functions offer three different idiomatic patterns for environment variable access in C++:
\begin{itemize}
  \item \textbf{Fail-fast:} \lstinline|getenv_or_exit|
  \item \textbf{Fallback:} \lstinline|getenv_or_default|
  \item \textbf{Exception-based:} \lstinline|getenv_or_throw|
\end{itemize}

They help enforce clear, consistent handling of environment variables and avoid scattered manual null checks throughout your codebase.

\end{document}
